\chapter{欲しい・たい — Deseos}
\unidadheader{22}{欲しい・たい}{Deseos}{Expresar lo que quieres hacer o tener}

\begin{tcolorbox}[vocabbox, title={📌 Vocabulario Nuevo}]
\begin{tabularx}{\linewidth}{llX}
\toprule
\textbf{Japonés} & \textbf{Español} & \textbf{Tipo} \\
\midrule
\vocabitem{欲しい}{ほしい}{querer (tener algo)}{adj-い}
\vocabitem{夢}{ゆめ}{sueño (aspiración)}{sust.}
\vocabitem{目標}{もくひょう}{meta / objetivo}{sust.}
\vocabitem{将来}{しょうらい}{futuro}{sust.}
\vocabitem{なりたい}{—}{quiero ser/convertirme en}{expr.}
\vocabitem{〜に行きたい}{—}{quiero ir a ~}{expr.}
\vocabitem{ぜひ}{—}{definitivamente / con gusto}{adv.}
\vocabitem{いつか}{—}{algún día}{adv.}
\vocabitem{できれば}{—}{si es posible}{expr.}
\vocabitem{どうしても}{—}{a toda costa / inevitablemente}{adv.}
\bottomrule
\end{tabularx}
\end{tcolorbox}

\subsection*{🗣️ Frases Clave}

\frase{\ruby{日本}{にほん}に\ruby{行}{い}きたいです。}{Quiero ir a Japón.}
\frase{あたらしいパソコンが\ruby{欲}{ほ}しいです。}{Quiero una computadora nueva.}
\frase{\ruby{将来}{しょうらい}、\ruby{医者}{いしゃ}になりたいです。}
  {En el futuro, quiero ser médico/a.}
\frase{いつかN1に\ruby{合格}{ごうかく}したいです。}
  {Algún día quiero aprobar el N1.}

\subsection*{⚙️ Gramática}

\patron{Querer hacer algo: たい}
  {V(ます-stem) + たいです / たくないです}
  {Solo se usa para deseos PROPIOS, no de terceros.
  \ej{\ruby{寿司}{すし}を\ruby{食}{た}べたいです。}{Quiero comer sushi.}
  \ej{もう\ruby{働}{はたら}きたくないです。}
    {Ya no quiero trabajar más.}
  \ej{どこに\ruby{行}{い}きたいですか？}{¿A dónde quieres ir?}}

\patron{Querer tener algo: 欲しい}
  {モノ + が + 欲しいです / 欲しくないです}
  {\ej{あたらしい\ruby{車}{くるま}が\ruby{欲}{ほ}しいです。}
    {Quiero un coche nuevo.}
  \ej{お\ruby{金}{かね}より\ruby{時間}{じかん}が\ruby{欲}{ほ}しいです。}
    {Quiero tiempo más que dinero.}}

\patron{Querer que alguien haga algo: てほしい}
  {V-て + ほしいです}
  {\ej{もっとゆっくり\ruby{話}{はな}してほしいです。}
    {Quiero que hables más despacio.}}

\begin{tcolorbox}[ejerciciobox, title={📝 Ejercicios Rápidos}]
\begin{enumerate}
  \item Di tres cosas que quieres hacer este año usando たいです.
  \item ¿Qué quieres tener? Usa 欲しいです.
  \item ¿Cuándo NO se puede usar たい? (pista: tercera persona)
\end{enumerate}
\vspace{0.4em}
\textbf{Respuestas:}
1) (personal)\quad
2) (personal)\quad
3) Para hablar de los deseos de terceros se usa 〜たがっています (está queriendo hacer ~) en lugar de たい.
\end{tcolorbox}

\begin{tcolorbox}[kanjibox, title={🀄 Kanjis de la Unidad}]
\begin{tabularx}{\linewidth}{cllXX}
\toprule
\textbf{Kanji} & \textbf{On} & \textbf{Kun} & \textbf{Significado} & \textbf{Vocab} \\
\midrule
\kanjirow{欲}{よく}{ほ}{desear / codicia}{\ruby{欲しい}{ほしい} / \ruby{欲求}{よっきゅう}}
\kanjirow{夢}{む}{ゆめ}{sueño}{\ruby{夢}{ゆめ} / \ruby{夢中}{むちゅう}}
\kanjirow{目}{もく}{め}{ojo / meta}{\ruby{目標}{もくひょう} / \ruby{目的}{もくてき}}
\kanjirow{標}{ひょう}{しるべ}{señal / meta}{\ruby{目標}{もくひょう} / \ruby{標準}{ひょうじゅん}}
\kanjirow{将}{しょう}{まさ}{futuro / general}{\ruby{将来}{しょうらい} / \ruby{将棋}{しょうぎ}}
\bottomrule
\end{tabularx}
\end{tcolorbox}
